\section{Source-clicks and tree orientations}

In this section, we study source-to-sink transformations from the point of view of admissible cuts. We first consider an arbitrary finite poset $P$ and a minimal element $s$, and give a criterion for reversing the covering relations incident with $s$ to be realized by the construction of Theorem \ref{generalizzazione teomio}. Whenever this criterion is satisfied, the poset obtained by the source-click is universally derived equivalent to $P$.

Source-to-sink transformations are classical in the representation theory of quivers, where they arise in the Bernstein--Gelfand--Ponomarev reflection functors \cite{BGP}. Their relation with derived equivalences was further developed by Happel \cite{Happel}. 

We then specialize to orientations of finite trees. In this case, the criterion is automatically satisfied at every source, so each source-click is realized by an admissible cut. We use the classical fact that any orientation of a finite tree can be reached from any other by a finite sequence of such clicks, and include a self-contained proof for completeness. This recovers Ladkani's result \cite[Corollary 1.8]{ladkani} that all orientations of a finite tree are universally derived equivalent. Related orientation-independence results have also been obtained in the stronger framework of strong stable equivalence; see Groth and Šťovíček \cite[Corollary 9.23]{Groth}.

Finally, we discuss why the analogous transitivity statement fails in general for graphs with cycles.


\subsection{An intrinsic criterion for source-clicks} \label{An intrinsic criterion for source-clicks}

Let $P$ be a finite poset and let $s$ be a minimal element of $P$. We denote by $\operatorname{Cov}_P(s)$ the set of elements covering $s$, that is, \[\operatorname{Cov}_P(s) \coloneq \{m \in P \colon s <_Pm \text{ and there is no } x \in P \text{ with } s<_Px<_Pm\}.\]

Since $s$ is a source in the oriented Hasse diagram of $P$, reversing all the covering relations incident with $s$ does not create an oriented cycle: after the reversal, $s$ becomes a sink. We denote by $\operatorname{cl}_s(P)$ the poset obtained by reversing every covering relation \[s <_Pm, \,\,\,\, m \in \operatorname{{Cov}_P(s)},\]
and taking the transitive closure. We call this operation a source-click at $s$.

The source-to-sink equivalence is the case $X=\{*\}$ of Ladkani's construction \cite[Corollary 1.7]{ladkani}. The point of the following proposition is to formulate this construction intrinsically for a given poset $P$: for a minimal element $s$, the relevant subset is necessarily $\operatorname{Cov}_P(s)$, and the admissible-cut condition gives a recognition criterion entirely in terms of the order structure of $P$.

\begin{proposition}%[General source-click criterion]
\label{General source-click criterion}
Let $P$ be a finite poset, let $s$ be a minimal element of $P$, and set
\[A=\{s\}, \,\,\,\, B=P \setminus\{s\}.\] Then \[M(s)=\operatorname{Cov}_P(s).\]
Moreover, the following conditions are equivalent:
\begin{enumerate}
    \item the partition $P=A \sqcup B$ is an admissible cut;
    \item for every pair of distinct elements $m,m' \in \operatorname{Cov_P(s)},$
    \[[m,\bullet]_B \cap [m', \bullet]_B= \emptyset\] and \[[\bullet, m]_B \cap [\bullet, m']_B = \emptyset.\]
\end{enumerate}
    Whenever these equivalent conditions hold, the associated negative poset coincides with the source-click of $P$ at $s$:
    \[P^-=\operatorname{cl}_s(P).\]
    In particular, $P$ and $\operatorname{cl}_s(P)$ are universally derived equivalent.
    
\end{proposition}

\begin{proof}
   Since $s$ is minimal, $A=\{s\}$ is an ideal of $P$, and therefore $B=P \setminus\{s\}$ is a filter.

   We first identify $M(s)$. By definition, \[M(s)=\operatorname{Min}_B\{b \in B \colon s \leq_Pb\}.\] An element $b \in B$ is minimal among the elements above $s$ if and only if there is no element $x \in P$ such that $s<_Px<_Pb$. This is equivalent to saying that $b$ covers $s$. Hence \[M(s)= \operatorname{Cov}_P(s).\]
   The ideal-filter condition is therefore automatic. Moreover, the compatibility condition is vacuous, since $A$ consists of a single element. Consequently, the partition $P=A \sqcup B$ is admissible if and only if the separation condition holds for the distinct elements of \[M(s)=\operatorname{Cov}_P(s),\] which proves the equivalence between $(i)$ and $(ii)$. 

   Assume now that these equivalent conditions hold. By Corollary \ref{The associated universally derived equivalent poset}, the mixed relations in the associated negative poset are
   \[b\leq_-s \iff \text{ there exists } m \in \operatorname{Cov}_P(s) \text{ such that } b \leq_B m.\]
   On the other hand, in the source-click $\operatorname{cl}_s(P)$, every edge $s<_Pm$ is replaced by $m<s$. Thus, for $b \in B$, there is an increasing path from $b$ to $s$ in $\operatorname{cl}_s(P)$ if and only if there exists $m \in \operatorname{Cov}_P(s)$ such that $b \leq_Bm$. Therefore 
   \[b \leq_{\operatorname{cl}_s(P)}s \iff b \leq_-s.\] The orders induced on $B$ are unchanged in both posets. Indeed, since $B$ is a filter of $P$, every chain in $P$ between elements of $B$ lies entirely in $B$, while after the source-click $s$ is a sink and therefore cannot create new relations between elements of $B$. Moreover, neither poset has a mixed relation of the form $s \leq b$. Hence $P^-= \operatorname{cl}_s(P)$. The universal derived equivalence follows from Corollary \ref{The associated universally derived equivalent poset}.

   
\end{proof}

\subsection{Tree orientations}

\textbf{Setup and notation.} Let $T$ be a finite tree with vertex set $V$ and edge set $E$. An orientation of $T$ is a partial order $P$ on $V$ whose covering relation, viewed as an undirected graph, equals $T$. Equivalently, an orientation is obtained by directing each edge of $T$ and taking the transitive closure. Since $T$ is acyclic, this always gives a partial order. We denote by $\operatorname{Or}(T)$ the set of orientations of $T$. A vertex $s$ is called a source of $P$ if it is minimal in $P$, i.e., for every edge joining $s$ to a vertex $v$, it holds that $s <_Pv$. Dually, $s$ is a sink if it is maximal in $P$. 

    Notice that, given $P \in \operatorname{Or(T)}$ and $x \neq y \in V$, $x \leq_Py$ if and only if the order increases at every step along the unique path from $x$ to $y$. 

For $P \in \operatorname{Or(T)}$ and a source $s$ of $P$, the source-click $cl_s(P)$ defined in Section \ref{An intrinsic criterion for source-clicks} amounts to reversing all edges incident to $s$; we also refer to this operation as a click at $s$.

\begin{corollary}%[Source step for tree orientations] 
\label{Source step}
Let $T$ be a finite tree, let $P \in \operatorname{Or(T)}$ with $|V| \geq 2$ and let $s$ be a source of $P$. Set $A=\{s\}$ and $B=V \setminus \{s\}$. Then $P=A \sqcup B$ is an admissible cut. Moreover, $M(s)=N(s)$, where $N(s)$ is the set of neighbours of $s$ in $T$, and the associated negative poset is the source-click of $P$ at $s$: \[P^-=\operatorname{cl}_s(P).\] In particular, $P$ and $cl_s(P)$ are universally derived equivalent.   
\end{corollary}

\begin{proof}
    Since the covering relation of $P$, viewed as an undirected graph, is $T$, and $s$ is a source, the elements covering $s$ are exactly its neighbours in $T$. Hence
    \[\operatorname{Cov}_P(s)=N(s).\] 
    Let $m \neq m'$ be two neighbours of $s$. The vertices $m$ and $m'$ belong to distinct connected components of $T \setminus \{s\}$. 

    Any element of $B$ which is comparable with $m$ through an increasing path in the Hasse diagram belongs to the same component of $T \setminus \{s\}$ as $m$. The same holds for $m'$. Since the two components are disjoint, both the upper and the lower cones of $m$
 and $m'$ in $B$ are disjoint:
\[[m, \bullet]_B \cap [m', \bullet]_B= \emptyset,\]
\[[\bullet, m]_B \cap [\bullet,m']_B= \emptyset.\] The conclusion now follows from Proposition \ref{General source-click criterion}.
\end{proof}
For $P,Q \in \operatorname{Or}(T)$, we write $P \to Q$ if $Q$ is obtained from $P$ by clicking at a source; that is, if $Q=cl_s(P)$ for some source $s$ of $P$. We denote by $\to ^*$ the reflexive-transitive closure of this relation. Thus $P \to ^*Q$ means that there exists a finite sequence of orientations \[P=P_0 \to P_1 \to \dots \to P_r=Q,\] or equivalently that $Q$ can be reached from $P$ by finitely many source-clicks.
The following combinatorial fact is classical; see \cite{BGP}. We include a self-contained proof for completeness.

\begin{proposition}%[Transitivity of clicks on a tree]
\label{Transitivity of clicks on a tree}
    Let $T$ be a finite tree. For every $P,Q \in \operatorname{Or}(T)$, one has $P \to^*Q$.
\end{proposition}
\begin{proof}
    We proceed by induction on $n=|V(T)|$. The cases $n=1$ and $n=2$ are immediate. Assume $n \geq 3$. Fix a leaf $\ell$ of $T$ that has a unique neighbour $p$ and set $T'=T\setminus\ell$, which is a tree with $n-1$ vertices.

    Every orientation $P \in \operatorname{Or}(T)$ is determined by its restriction $P' \in \operatorname{Or}(T')$, together with the direction of the edge $\ell p$. We say that the leaf is in the state $L$ if $\ell <_Pp$, and in state $H$ if $p<_P\ell$. We write $P=(P',L)$ or $P=(P',H)$.

    Interactions of clicks with this decomposition:
    \begin{itemize}
        \item $\ell$ is in state $L$ if and only if $\ell$ is a source of $P$. In this case, clicking $\ell$ changes the state from $L$ to $H$ without changing the restriction to $T'$.
        \item A vertex $v\in V(T') \setminus \{p\}$ is a source of $P$ if and only if it is a source of $P'$. (Because its edges coincide in $T$ and $T'$.) In this case, clicking $v$ induces on $T'$ the same operation as clicking $v$ in $P'$, and it leaves the leaf state unchanged.
        \item The vertex $p$ is a source of $P$ if and only if it is a source of $P'$ and the leaf is in the state $H$. In this case, clicking $p$ induces the click at $p$ on $T'$ and changes the leaf state from $H$ to $L$.
    \end{itemize}

By the induction hypothesis, any two orientations of $T'$ can be connected by a sequence of source-clicks. Let $\pi$ be such a sequence starting from $P'$. We lift $\pi$ to $T$ by performing its clicks in the same order and inserting, when necessary, clicks at the leaf $\ell$, which leave the restriction to $T'$ unchanged. We read $\pi$ from left to right. 

A click at vertex $v \neq p$ is lifted to the same click at $v$ in $T$. 

A click at $p$ is lifted as follows: if the current leaf state is $H$, click $p$; if the current leaf state is $L$, first click $\ell$, changing the state to $H$, and then click $p$. In this way the restriction to $T'$ follows exactly the sequence $\pi$.

Moreover, after every lifted click at $p$, the leaf state is $L$. Clicks at vertices different from $p$ and $\ell$ do not change the leaf state. Thus, after lifting $\pi$, the final restriction to $T'$ is the desired one, and the final state is determined explicitly.

Now let $Q=(Q',\eta) \in \operatorname{Or}(T)$ be arbitrary, with state $\eta \in \{L,H\}$. By induction, choose a source-click sequence in $T'$ from $P'$ to $Q'$, and lift it to $T$ as above. This gives an orientation of the form $(Q', \epsilon)$. If $\epsilon = \eta$, we are done. If $\epsilon=L$ and $\eta=H$, click the leaf $\ell$, which changes only the leaf state.

It remains only to see that the opposite change, from $H$ to $L$ at fixed $Q'$, is also possible. Choose an orientation $R' \in \operatorname{Or}(T')$ in which $p$ is a source, for example by orienting every edge away from $p$. By the induction hypothesis, there is a source-click sequence from $Q'$ to $R'$, and another one from $cl_p(R')$ back to $Q'$. Concatenating these sequences with the click at $p$ gives a closed source-click walk based at $Q'$, that is, a sequence of source-clicks on $T'$ that starts and ends at $Q'$. This walk contains a click at $p$. Lifting this closed walk from the state $(Q',H)$ returns the restriction to $Q'$, and after the last lifted click at $p$, the leaf state is $L$. Hence $(Q',H) \to^*(Q',L)$. 

Therefore both leaf states are reachable over the same restriction $Q'$, and so $P \to ^*Q$. This completes the induction.
\end{proof}
\begin{theorem} %[tree-exhaustiveness]
\label{tree-exhaustiveness}
    Let $T$ be a finite tree. Starting from any orientation $P \in \operatorname{Or}(T)$, every orientation of $T$ can be reached by a finite sequence of applications of the general construction of Theorem \ref{generalizzazione teomio} to single-source decompositions. In particular, all orientations of $T$ are universally derived equivalent.
\end{theorem}
\begin{proof}
    Let $P,Q \in \operatorname{Or}(T).$ By Proposition \ref{Transitivity of clicks on a tree}, there exists a finite sequence of source-clicks from $P$ to $Q$. By Corollary \ref{Source step}, each source-click is realized by the construction of Theorem \ref{generalizzazione teomio} applied to a decomposition of the form $A=\{s\}$, where $s$ is a source of the current orientation. Hence consecutive orientations in the sequence are universally derived equivalent. Since universal derived equivalence is transitive, $P$ and $Q$ are universally derived equivalent. Therefore all orientations of $T$ are universally derived equivalent.
\end{proof}

\begin{remark}
    If $s$ is a source with at least two neighbours, then $M(s)=N(s)$ has more than one element. Thus the source step of Corollary \ref{Source step} cannot, in general, be obtained from the singleton case of Corollary \ref{Reconstructing a poset from a cut}. Already in the path $A_3$, the middle vertex can be the unique source, and clicking it requires the general construction.
\end{remark}

\begin{remark}
   The transitivity argument above uses that the underlying graph is a tree. For a graph with cycles, the analogous statement for acyclic orientations and graph-theoretic source-clicks can fail.

   Indeed, suppose that a graph contains a cycle $C$, and choose one of the two possible directions around $C$. For an acyclic orientation $O$, define $\sigma_C(O)$ as the number of edges of $C$ oriented in the chosen direction minus the number of edges oriented in the opposite direction. 

   A graph-theoretic source-click does not change $\sigma_C$. If the clicked vertex does not lie on $C$, this is clear. If it lies on $C$, then the two edges of $C$ incident with it are both reversed. Since the clicked vertex is a source, before the click one of these two edges is oriented in the chosen direction and the other one is oriented in the opposite direction. After the click, the two contributions are exchanged. Hence $\sigma_C$ is invariant under graph-theoretic source-clicks.

   On the other hand, acyclic orientations with different values of $\sigma_C$ exist. Indeed, write $C=(v_1, \dots, v_r,v_1)$ in the chosen direction and orient every edge according to a total order extending $v_1<\dots<v_r$. This gives $\sigma_C=r-2$, while reversing the total order gives $\sigma_C=2-r$.

   Therefore, for graphs with cycles, such clicks cannot in general connect all acyclic orientations. Notice that, unlike the tree case, a graph-theoretic source-click need not coincide with the poset source-click $\operatorname{cl_s}$ of Section \ref{An intrinsic criterion for source-clicks}: an edge of an acyclic orientation may become redundant after taking the transitive closure and hence not correspond to a covering relation. Thus $\sigma_C$ is an obstruction to transitivity for graph orientations, rather than in general an invariant of the poset source-click operation.

   In particular, the presence of a cycle does not by itself prevent an individual poset source-click from being realized by an admissible cut; this is governed by the criterion of Proposition \ref{General source-click criterion}.
\end{remark}